\subsection{Two waves of convergence}
The five maps show two distinct waves in the shrinking of Britain as seen from London. The first came on the roads. Between 1750 and the 1830s the fitted coach speed on a turnpike rose from {{mph.1750|f1}} to {{mph.1830|f1}} miles per hour of moving time, and the area within a day grew {{ratio.1830_1750|f1}}-fold. This is consistent with the timing of the gains that \citet{gerhold2014} and \citet{bogart2005} attribute to turnpikes, vehicles and horses, and it is a reminder that the convergence of the coaching age was large in its own right: in the model, the 1830s traveller reached in a day what had taken the 1750 traveller several days. The 1819 fit, at {{mph.1819|f1}} mph, sits between the two, as the literature's decade averages would predict.

The second wave came with the railway, and it was faster. Between the 1836 timetables and the 1851 network, about fifteen years, the area within a day grew {{ratio.1850_1830|f1}}-fold, and the shape of the reachable area changed. A coach-era isochrone is a compact, nearly circular patch; the railway isochrone of 1850 is ragged, with its inner hours drawn out along the trunk lines like the spokes of a star, while land between the lines remains far slower to reach. This is what one would expect of a network whose speed lay almost entirely on its links, and it matches the emphasis in the railway literature on the uneven way in which lines and stations reached the population \citep{gregory2010,bogart2022}. Between 1850 and 1910 the gain in area was smaller, {{ratio.1910_1850|f1}}-fold, but the network filled in, and by 1910 almost all of England and Wales, and Edinburgh and Glasgow, were within a day. The twentieth century then added the rest of the mainland, by road as much as by rail. The model's day is a blunt instrument for the period after 1910, because once most of the island is inside the line the area can no longer grow much; the six-hour area in Table~\ref{tab:area} shows that convergence continued long after the twelve-hour area saturated.

\subsection{Why calibrate}
Historical network models usually assume speeds. The exercise here shows why it is worth fitting them where timetables exist and testing them where they do not. With one parameter, the 1830s model reproduces {{fits.1830.n|int}} scheduled coach times (median error {{fits.1830.loo.mdape|f0}}\% leaving each out) and an independent list of {{valid.1830_mail1837.n|int}} mail coaches (median error {{valid.1830_mail1837.mdape|f0}}\%), which gives some confidence that the turnpike map, the assumption about non-turnpike roads and the walking links are adequate for the period. The same structure fitted to the much sparser 1750 evidence does not reach the standard set in advance: its leave-one-out median error is {{fits.1750.loo.mdape|f0}}\%, partly because the cited journeys disagree among themselves, partly because a duration in days does not say how many hours were spent on the road, and partly because the network of 1750 is coded only as turnpike or not, when the condition of an unturnpiked road could range from good to impassable. For the railway eras the comparison with the fastest documented trains exposes a different issue: an average train is not the fastest train. The model is slower than the documented fastest trains by a median {{valid.1850_all.mdape|f0}}\% around 1850 and {{valid.1910_all.mdape|f0}}\% around 1910, a gap as large as some of the differences between eras. Which is right depends on the question. A traveller who could choose their train, and pay for it, could go further than the maps show; a traveller who took what came would do worse, because the model has no timetable waits beyond its half-hour allowance.

\subsection{Limitations}
Several limitations follow from the data and the design. First, there are no timetables. Coaches and trains are assumed to be available when the traveller arrives, apart from a fixed wait, and connections, night services and Sunday working are ignored. Second, the 1910 network is a reconstruction. Lines that closed after 1910 are inferred as straight links between closed stations; they are shorter than the real lines, may join stations that were never connected, and miss closed lines whose stations were within 2 km of a surviving or 1851 line. Lines opened after 1910 and still open are not removed. The sensitivity check without the inferred lines shows the scale of their effect. Third, the main-line classification for 1850 and 1910 uses today's main lines as a proxy. Fourth, Scotland has no turnpike map, and its main roads are proxied by today's trunk and primary roads; this affects only regions beyond a day from London in the coach eras, but it would matter for a study centred in Scotland. Fifth, the 1750 calibration rests on twelve journeys, several from narrative sources quoting advertisements, which promise rather than record. Sixth, the twelve-hour day is a choice; Table~\ref{tab:sens} shows that an eight or sixteen-hour day changes the areas more than any other assumption, though not their order. Seventh, the model ignores terrain, weather and season, fares and the cost of a horse, and assumes every traveller could ride in a coach or train at all, which in 1750 most people could not afford. Finally, the town list and the origin are fixed: a day from Manchester or Bristol would look different, and the same code could produce it.

\subsection{Future work}
Three sources would remove the largest uncertainties. Leunig's transcriptions of railway timetables for important and minor routes in 1850, 1870, 1887 and 1910, and for 48 routes between 1910 and 2008, are deposited with the UK Data Service and would allow train speeds to be fitted and tested route by route. The full time-dynamic rail lines file would replace the inferred lines of 1910 with the real network. And a systematic collection of dated coach advertisements for the 1740s to 1760s, of the kind Gerhold assembled, would put the 1750 map on the same footing as the 1830s. With timetables, the model could also be extended from a fastest-path model to one that respects departure times, which would bring it closer to what a traveller actually experienced.
